\subsection{Rolle von control.py}
\label{subsec:rolle-control-py}

\path{tools/control.py} bündelt die öffentlichen Arbeitsabläufe. Seine Befehle
decken Initialisierung, Diagnose, Installation, Entwicklung, Konfiguration,
Datenbank, Quality, Tests, Builds, Container, Tauri, Versionierung, Release und
Dokumentation ab. Hilfe und Gruppenansichten strukturieren diese Oberfläche
\parencite{kleiveist2026control,kleiveist2026quality-implementation}.
Zu den wichtigsten Einstiegen gehören \texttt{init}, \texttt{doctor},
\texttt{install}, \texttt{run}, \texttt{quality}, \texttt{test} und
\texttt{build}. Die Gruppen \texttt{db}, \texttt{tauri},
\texttt{container} und \texttt{release} ergänzen nur die jeweils benötigten
Fachabläufe.

\begin{figure}[H]
  \centering
  \resizebox{0.98\textwidth}{!}{%
  \begin{tikzpicture}
    \node[casePrimary,minimum width=52mm,minimum height=14mm] (control) at (0,0)
      {\texttt{control.py}\\gemeinsamer öffentlicher Einstieg};

    \node[caseBox,minimum width=37mm,minimum height=25mm] (setup) at (-6.3,-2.4)
      {\textbf{Setup / Diagnose}\\\texttt{init / doctor / install}};
    \node[caseBox,minimum width=37mm,minimum height=25mm] (development) at (-2.1,-2.4)
      {\textbf{Entwicklung / Qualität}\\\texttt{run / stop / quality}\\\texttt{test / build / console}};
    \node[caseBox,minimum width=37mm,minimum height=25mm] (subsystems) at (2.1,-2.4)
      {\textbf{Subsysteme}\\\texttt{config / db}\\\texttt{tauri / container}};
    \node[caseBox,minimum width=37mm,minimum height=25mm] (delivery) at (6.3,-2.4)
      {\textbf{Auslieferung / Pflege}\\\texttt{version / release / docs}};

    \coordinate (branch) at (0,-0.95);
    \draw[caseFlow] (control.south) -- (branch) -| (setup.north);
    \draw[caseFlow] (branch) -| (development.north);
    \draw[caseFlow] (branch) -| (subsystems.north);
    \draw[caseFlow] (branch) -| (delivery.north);
  \end{tikzpicture}}
  \caption{Befehlsgruppen des zentralen Projekt-Toolings}
  \label{fig:control-command-map}
  \smallskip
  {\footnotesize Quelle: Eigene Darstellung auf Grundlage von
  \textcite{kleiveist2026control}, \textcite{kleiveist2026tooling} und
  \textcite{kleiveist2026quality-implementation}.\par}
\end{figure}


Der Governance-Commit teilte den zuvor 592 physische Zeilen umfassenden
Dispatcher auf. \path{tools/control.py} umfasst im untersuchten Stand 197
Zeilen und fungiert als Composition Root; Parserdefinitionen liegen in
\path{tools/control_parser.py}, die eigentliche Prüfung im modularen Paket
\path{tools/quality/}. Eine Rückabhängigkeit dieses Pakets auf den Dispatcher
besteht nicht. Die Bewertung folgt damit Verantwortung, Kopplung und
Delegation und nicht allein der Dateilänge
\parencite{kleiveist2026governance-commit}.

Ohne Aktion entspricht \texttt{quality} dem Unterbefehl \texttt{check}.
Fokussierte Aktionen sind \texttt{size}, \texttt{complexity},
\texttt{architecture}, \texttt{lint} und \texttt{format}; als Optionen stehen
\texttt{--config} und \texttt{--format text|json} bereit. Jede Aktion lädt die
Policy, scannt das Repository und validiert Exceptions. Die vollständige
Aktion verbindet zusätzlich Größen- und Komplexitätsregeln,
Architekturprüfung, Sprachwerkzeuge, Compiler- und Formatchecks.

Exitcode~0 setzt erfolgreiche Toolchecks voraus; aktive Fehlerbefunde dürfen
nicht vorliegen. Ein nicht unterdrückter \emph{ERROR}-Befund oder ein
fehlgeschlagenes erforderliches Werkzeug liefert~1. Parser-, Konfigurations-
und Scanfehler liefern~2, ein Tastaturabbruch auf Dispatcher-Ebene~130. Die
Stufen \emph{INFO}, \emph{WARNING} und die starke Warnung blockieren allein
nicht. \emph{INFO} wird im Modell und JSON-Report gezählt, aber von keiner
v1-Regel erzeugt. Externe Lint-, Compiler- und Formatfehler erscheinen als
fehlgeschlagene Checks mit Werkzeugausgabe, nicht immer als strukturiertes
Finding mit CQ-ID.

Tabelle~\ref{tab:quality-rule-groups} verdichtet die 13 implementierten
Regelkennungen. Die CQ-Severity hängt von der jeweiligen Schwelle ab; AR- und
EX-Befunde sind Fehler.

\begin{table}[H]
  \centering
  \scriptsize
  \renewcommand{\arraystretch}{1.08}
  \caption{Regelgruppen der Governance~v1}
  \label{tab:quality-rule-groups}
  \begin{tabularx}{\textwidth}{@{}p{0.15\textwidth}p{0.31\textwidth}>{\raggedright\arraybackslash}X p{0.16\textwidth}@{}}
    \toprule
    \textbf{Gruppe} & \textbf{Regel-IDs} & \textbf{Prüfgegenstand} &
      \textbf{Severity} \\
    \midrule
    Codequalität & \texttt{CQ001}, \texttt{CQ002},
      \texttt{CQ101}--\texttt{CQ104}, \texttt{CQ201} & Datei-, Funktions-
      und Klassengröße, Komplexität, Verschachtelung und Parameter &
      schwellenabhängig \\
    Architektur & \texttt{AR001}--\texttt{AR004} & Schicht- und
      Frameworkabhängigkeiten, Feature-Interna und Router-Heuristik & ERROR \\
    Exceptions & \texttt{EX001}, \texttt{EX002} & ungültige oder abgelaufene
      Ausnahmeeinträge & ERROR \\
    \bottomrule
  \end{tabularx}
  \smallskip
  {\footnotesize Quelle: Eigene Darstellung nach
  \textcite{kleiveist2026quality-implementation}.\par}
\end{table}

Repository-Findings enthalten Pfad, optionale Zeile und Symbol, ID, Name,
Severity, Istwert, Schwelle, Detail und Handlungsempfehlung. Unterdrückte
Befunde bleiben einschließlich Begründung sichtbar; die Summary zählt
Severity-Stufen und Exceptions. Der Textbericht weist INFO nicht gesondert
aus, der JSON-Bericht schon. Außerdem stammt die Zahl \texttt{files\_checked}
praktisch nur aus der Größenprüfung, weshalb fokussierte Komplexitäts- und
Architekturläufe trotz geprüfter Quellen null Dateien melden können. Diese
Reporting-Grenze verändert den Befundstatus nicht.
\beginnerinput{workspace/statement/500_tooling_workflow/510}
